% =====================================================================
%  Cadenas de idioma: Español
%  Seleccionado por \booklang=es (ver main.tex). Define el idioma de
%  babel, los nombres de los entornos de teoremas, los títulos de las
%  partes y el subtítulo del libro para que una sola fuente compile en
%  cualquiera de los dos idiomas.
% =====================================================================
\newcommand{\babellang}{spanish}

% ----- Nombres de entornos de teoremas ------------------------------
\newcommand{\nameDefinition}{Definici\'on}
\newcommand{\nameExample}{Ejemplo}
\newcommand{\nameRemark}{Observaci\'on}
\newcommand{\nameTheorem}{Teorema}
\newcommand{\nameLemma}{Lema}
\newcommand{\nameProposition}{Proposici\'on}
\newcommand{\nameCorollary}{Corolario}
\newcommand{\nameHypothesis}{Hip\'otesis}
\newcommand{\nameConjecture}{Conjetura}
\newcommand{\nameApplication}{Aplicaci\'on}

% ----- Títulos de glosario / listas ---------------------------------
\newcommand{\glossaryTitle}{Glosario}
\newcommand{\acronymTitle}{Acr\'onimos}

% ----- Títulos de las partes ----------------------------------------
\newcommand{\partFoundations}{Fundamentos}
\newcommand{\partFiveGs}{Las Cinco Ges}
\newcommand{\partPractice}{Pr\'actica}

% ----- Portada -------------------------------------------------------
% Título completo (igual que el TFG): «Geometric Deep Learning, una
% perspectiva matemática» — se compone de título + subtítulo en la portada.
% La línea de las 5 Ges se conserva como lema.
\newcommand{\booktitle}{Geometric Deep Learning}
\newcommand{\booksubtitle}{una perspectiva matem\'atica}
\newcommand{\booktagline}{Rejillas, Grupos, Grafos, Geod\'esicas y Gauges}
\newcommand{\bookauthor}{Jorge Mu\~noz Fuentes}
\newcommand{\bookpublisher}{Fractal Manifold S.L.}

% ----- Texto de contraportada (usado por la portada envolvente) ------
% Sinopsis breve condensada del resumen, ajustada a una contraportada.
\newcommand{\bookblurb}{%
  El aprendizaje profundo geom\'etrico es un marco unificador que deriva
  las arquitecturas de redes neuronales a partir de las simetr\'ias y la
  estructura geom\'etrica de los datos sobre los que operan. Partiendo de
  la invarianza y la equivarianza, este libro construye la teor\'ia a
  trav\'es de las cinco Ges --- rejillas, grupos, grafos, geod\'esicas y
  gauges --- y la conecta con los modelos que definen el aprendizaje
  autom\'atico moderno: redes convolucionales, redes equivariantes a
  grupos, redes neuronales de grafos y m\'as all\'a.}
